% --- Archon physics formalization source begin ---
% archon:physics
% archon:covers PhyXMiniProblems/problem_phyx_mini_0846.lean
% archon:source-report reports/phyx_mini/problem_phyx_mini_0846.source.json

\chapter{Physics problem phyx\_mini\_0846}
\label{ch:PhyXMiniProblems_problem_phyx_mini_0846}

\paragraph{Problem source.}
\#\# Question

What is the net electric flux through the cylinder?

\#\# Answer choices

- A: A: -1.2N m\textasciicircum{}2/C

- B: B: -10N m\textasciicircum{}2/C

- C: C: -2N m\textasciicircum{}2/C

- D: D: 0N m\textasciicircum{}2/C

\#\# Figure caption (auxiliary; use the image as primary evidence)

The image depicts a cylindrical shape with a uniform electric field, represented by pink arrows, passing horizontally from left to right across and through the cylinder. The arrows indicate the direction of the electric field (\textbackslash{}( \textbackslash{}vec\{E\} \textbackslash{})).

Key elements include:

- **Cylinder**: A 3D cylindrical shape is shown with one end visible and the other end partially shown with a dashed line.
- **Electric Field**: Pink arrows representing the electric field (\textbackslash{}( \textbackslash{}vec\{E\} \textbackslash{})) are directed parallel to the axis of the cylinder. Some arrows are also shown entering and exiting the cylinder, indicating that the field uniformly permeates the space.
- **Dimension Label**: The cylinder’s diameter is labeled as \textbackslash{}( 2R \textbackslash{}) with a vertical arrow inside the cylinder showing this measurement.
- **Text**: The electric field vector symbol (\textbackslash{}( \textbackslash{}vec\{E\} \textbackslash{})) is written next to the arrows just outside the left side of the cylinder.

The image is likely used to illustrate concepts related to electric fields and cylinders, possibly in a physics context involving Gauss's Law or field distribution around conductors.

\#\# Dataset metadata

Electrostatics; Physical Model Grounding Reasoning, Multi-Formula Reasoning

\paragraph{Recorded answer/context.}
D: D: 0N m\textasciicircum{}2/C

\paragraph{Figure/image path.}
/root/proposal\_for\_physic/science-mango/phyx\_mini\_run/phyx\_data/test\_image/846.png

\paragraph{Formalization target.}
create a compiling Lean file with sorry bodies at `PhyXMiniProblems/problem\_phyx\_mini\_0846.lean`. The Lean declarations must preserve the physical quantities, dimensions or dimensional roles, figure labels, governing-law hypotheses, and final relation expressed by this problem.
Use Mathlib/Physlib names found through LeanExplore where available. If a physics API is missing, introduce faithful local abstractions rather than scalar placeholder aliases.

\begin{theorem}[Physics formalization target]
\label{thm:physics:phyx_mini_0846:target}
\lean{PhyXMiniProblems.ProblemPhyXMini0846.netElectricFluxThroughCylinder_eq_zero}
\uses{def:physics:phyx-mini-0846:phyxminiproblems-problemphyxmini0846-electricfluxinnewtonsquaremeterspercoulomb, def:physics:phyx-mini-0846:phyxminiproblems-problemphyxmini0846-uniformaxialfieldcylindersetup, def:physics:phyx-mini-0846:phyxminiproblems-problemphyxmini0846-matchessuppliedcylinderfigure, def:physics:phyx-mini-0846:phyxminiproblems-problemphyxmini0846-hasphysicalclosedcylindergeometry, def:physics:phyx-mini-0846:phyxminiproblems-problemphyxmini0846-isuniformelectricfieldparalleltocylinderaxis, def:physics:phyx-mini-0846:phyxminiproblems-problemphyxmini0846-satisfiesclosedcylinderelectricfluxlaws}
The net electric flux through the closed cylinder is zero.
\end{theorem}
\begin{proof}
The uniform field is parallel to the cylinder axis.  The orientation factors of the left cap, curved wall, and right cap are respectively $-1$, $0$, and $1$.  The flux law therefore gives contributions $-EA$, $0$, and $EA$.  Substitute these into finite additivity of the three surface contributions; they sum to zero.
\end{proof}

% --- Archon declaration topology begin ---
\section{Declaration topology}
\begin{definition}[electric Field Strength Dimension]
  \label{def:physics:phyx-mini-0846:phyxminiproblems-problemphyxmini0846-electricfieldstrengthdimension}
  \lean{PhyXMiniProblems.ProblemPhyXMini0846.electricFieldStrengthDimension}
  The dimension M L T⁻² C⁻¹ of electric-field strength
\end{definition}
\begin{proof}
  This is the declared model definition of electric Field Strength Dimension.
\end{proof}
\begin{definition}[electric Flux Dimension]
  \label{def:physics:phyx-mini-0846:phyxminiproblems-problemphyxmini0846-electricfluxdimension}
  \lean{PhyXMiniProblems.ProblemPhyXMini0846.electricFluxDimension}
  The dimension M L³ T⁻² C⁻¹ of electric flux, i.e. N m² / C
\end{definition}
\begin{proof}
  This is the declared model definition of electric Flux Dimension.
\end{proof}
\begin{definition}[Length Quantity]
  \label{def:physics:phyx-mini-0846:phyxminiproblems-problemphyxmini0846-lengthquantity}
  \lean{PhyXMiniProblems.ProblemPhyXMini0846.LengthQuantity}
  A nonnegative, unit-independent physical length
\end{definition}
\begin{proof}
  This is the declared model definition of Length Quantity.
\end{proof}
\begin{definition}[Area Quantity]
  \label{def:physics:phyx-mini-0846:phyxminiproblems-problemphyxmini0846-areaquantity}
  \lean{PhyXMiniProblems.ProblemPhyXMini0846.AreaQuantity}
  A nonnegative, unit-independent physical area
\end{definition}
\begin{proof}
  This is the declared model definition of Area Quantity.
\end{proof}
\begin{definition}[Electric Field Strength Quantity]
  \label{def:physics:phyx-mini-0846:phyxminiproblems-problemphyxmini0846-electricfieldstrengthquantity}
  \lean{PhyXMiniProblems.ProblemPhyXMini0846.ElectricFieldStrengthQuantity}
  \uses{def:physics:phyx-mini-0846:phyxminiproblems-problemphyxmini0846-electricfieldstrengthdimension}
  A nonnegative magnitude of electric-field strength
\end{definition}
\begin{proof}
  This is the declared model definition of Electric Field Strength Quantity.
\end{proof}
\begin{definition}[Electric Flux Quantity]
  \label{def:physics:phyx-mini-0846:phyxminiproblems-problemphyxmini0846-electricfluxquantity}
  \lean{PhyXMiniProblems.ProblemPhyXMini0846.ElectricFluxQuantity}
  \uses{def:physics:phyx-mini-0846:phyxminiproblems-problemphyxmini0846-electricfluxdimension}
  A signed electric flux through an oriented surface
\end{definition}
\begin{proof}
  This is the declared model definition of Electric Flux Quantity.
\end{proof}
\begin{definition}[length In Meters]
  \label{def:physics:phyx-mini-0846:phyxminiproblems-problemphyxmini0846-lengthinmeters}
  \lean{PhyXMiniProblems.ProblemPhyXMini0846.lengthInMeters}
  \uses{def:physics:phyx-mini-0846:phyxminiproblems-problemphyxmini0846-lengthquantity}
  SI readout of a physical length, in metres
\end{definition}
\begin{proof}
  This is the declared model definition of length In Meters.
\end{proof}
\begin{definition}[area In Square Meters]
  \label{def:physics:phyx-mini-0846:phyxminiproblems-problemphyxmini0846-areainsquaremeters}
  \lean{PhyXMiniProblems.ProblemPhyXMini0846.areaInSquareMeters}
  \uses{def:physics:phyx-mini-0846:phyxminiproblems-problemphyxmini0846-areaquantity}
  SI readout of a physical area, in square metres
\end{definition}
\begin{proof}
  This is the declared model definition of area In Square Meters.
\end{proof}
\begin{definition}[electric Field Strength In Newtons Per Coulomb]
  \label{def:physics:phyx-mini-0846:phyxminiproblems-problemphyxmini0846-electricfieldstrengthinnewtonspercoulomb}
  \lean{PhyXMiniProblems.ProblemPhyXMini0846.electricFieldStrengthInNewtonsPerCoulomb}
  \uses{def:physics:phyx-mini-0846:phyxminiproblems-problemphyxmini0846-electricfieldstrengthquantity}
  SI readout of electric-field strength, in newtons per coulomb
\end{definition}
\begin{proof}
  This is the declared model definition of electric Field Strength In Newtons Per Coulomb.
\end{proof}
\begin{definition}[electric Flux In Newton Square Meters Per Coulomb]
  \label{def:physics:phyx-mini-0846:phyxminiproblems-problemphyxmini0846-electricfluxinnewtonsquaremeterspercoulomb}
  \lean{PhyXMiniProblems.ProblemPhyXMini0846.electricFluxInNewtonSquareMetersPerCoulomb}
  \uses{def:physics:phyx-mini-0846:phyxminiproblems-problemphyxmini0846-electricfluxquantity}
  SI readout of signed electric flux, in newton-square-metres per coulomb
\end{definition}
\begin{proof}
  This is the declared model definition of electric Flux In Newton Square Meters Per Coulomb.
\end{proof}
\begin{definition}[Cylinder Surface Part]
  \label{def:physics:phyx-mini-0846:phyxminiproblems-problemphyxmini0846-cylindersurfacepart}
  \lean{PhyXMiniProblems.ProblemPhyXMini0846.CylinderSurfacePart}
  The three pieces of the closed cylindrical Gaussian surface
\end{definition}
\begin{proof}
  This is the declared model definition of Cylinder Surface Part.
\end{proof}
\begin{definition}[Horizontal Direction]
  \label{def:physics:phyx-mini-0846:phyxminiproblems-problemphyxmini0846-horizontaldirection}
  \lean{PhyXMiniProblems.ProblemPhyXMini0846.HorizontalDirection}
  Horizontal directions needed to transcribe the side-view figure
\end{definition}
\begin{proof}
  This is the declared model definition of Horizontal Direction.
\end{proof}
\begin{definition}[Uniform Axial Field Cylinder Figure]
  \label{def:physics:phyx-mini-0846:phyxminiproblems-problemphyxmini0846-uniformaxialfieldcylinderfigure}
  \lean{PhyXMiniProblems.ProblemPhyXMini0846.UniformAxialFieldCylinderFigure}
  \uses{def:physics:phyx-mini-0846:phyxminiproblems-problemphyxmini0846-lengthquantity, def:physics:phyx-mini-0846:phyxminiproblems-problemphyxmini0846-cylindersurfacepart, def:physics:phyx-mini-0846:phyxminiproblems-problemphyxmini0846-horizontaldirection}
  Typed data read directly from the supplied raster figure. The diameter label is a physical length; its printed text and its relation to the radius are recorded separately below
\end{definition}
\begin{proof}
  This is the declared model definition of Uniform Axial Field Cylinder Figure.
\end{proof}
\begin{definition}[Closed Circular Cylinder]
  \label{def:physics:phyx-mini-0846:phyxminiproblems-problemphyxmini0846-closedcircularcylinder}
  \lean{PhyXMiniProblems.ProblemPhyXMini0846.ClosedCircularCylinder}
  \uses{def:physics:phyx-mini-0846:phyxminiproblems-problemphyxmini0846-lengthquantity, def:physics:phyx-mini-0846:phyxminiproblems-problemphyxmini0846-areaquantity}
  An abstract closed circular cylinder in three-dimensional physical space. cylinderAndBoundary identifies the spatial region on which the displayed field is asserted to be uniform; it is not a scalar substitute for the geometry
\end{definition}
\begin{proof}
  This is the declared model definition of Closed Circular Cylinder.
\end{proof}
\begin{definition}[Uniform Axial Field Cylinder Setup]
  \label{def:physics:phyx-mini-0846:phyxminiproblems-problemphyxmini0846-uniformaxialfieldcylindersetup}
  \lean{PhyXMiniProblems.ProblemPhyXMini0846.UniformAxialFieldCylinderSetup}
  \uses{def:physics:phyx-mini-0846:phyxminiproblems-problemphyxmini0846-electricfieldstrengthquantity, def:physics:phyx-mini-0846:phyxminiproblems-problemphyxmini0846-electricfluxquantity, def:physics:phyx-mini-0846:phyxminiproblems-problemphyxmini0846-cylindersurfacepart, def:physics:phyx-mini-0846:phyxminiproblems-problemphyxmini0846-uniformaxialfieldcylinderfigure, def:physics:phyx-mini-0846:phyxminiproblems-problemphyxmini0846-closedcircularcylinder}
  The physical setup and its independent flux observables. In particular, netElectricFlux is not defined from the displayed answer 0
\end{definition}
\begin{proof}
  This is the declared model definition of Uniform Axial Field Cylinder Setup.
\end{proof}
\begin{definition}[Matches Supplied Cylinder Figure]
  \label{def:physics:phyx-mini-0846:phyxminiproblems-problemphyxmini0846-matchessuppliedcylinderfigure}
  \lean{PhyXMiniProblems.ProblemPhyXMini0846.MatchesSuppliedCylinderFigure}
  \uses{def:physics:phyx-mini-0846:phyxminiproblems-problemphyxmini0846-lengthinmeters, def:physics:phyx-mini-0846:phyxminiproblems-problemphyxmini0846-uniformaxialfieldcylindersetup}
  Primary-image evidence, including the field direction and the 2R label
\end{definition}
\begin{proof}
  This is the declared model definition of Matches Supplied Cylinder Figure.
\end{proof}
\begin{definition}[Has Physical Closed Cylinder Geometry]
  \label{def:physics:phyx-mini-0846:phyxminiproblems-problemphyxmini0846-hasphysicalclosedcylindergeometry}
  \lean{PhyXMiniProblems.ProblemPhyXMini0846.HasPhysicalClosedCylinderGeometry}
  \uses{def:physics:phyx-mini-0846:phyxminiproblems-problemphyxmini0846-lengthinmeters, def:physics:phyx-mini-0846:phyxminiproblems-problemphyxmini0846-areainsquaremeters, def:physics:phyx-mini-0846:phyxminiproblems-problemphyxmini0846-uniformaxialfieldcylindersetup}
  Regularity and circular-end geometry of the physical cylinder
\end{definition}
\begin{proof}
  This is the declared model definition of Has Physical Closed Cylinder Geometry.
\end{proof}
\begin{definition}[Is Uniform Electric Field Parallel To Cylinder Axis]
  \label{def:physics:phyx-mini-0846:phyxminiproblems-problemphyxmini0846-isuniformelectricfieldparalleltocylinderaxis}
  \lean{PhyXMiniProblems.ProblemPhyXMini0846.IsUniformElectricFieldParallelToCylinderAxis}
  \uses{def:physics:phyx-mini-0846:phyxminiproblems-problemphyxmini0846-electricfieldstrengthinnewtonspercoulomb, def:physics:phyx-mini-0846:phyxminiproblems-problemphyxmini0846-uniformaxialfieldcylindersetup}
  The displayed field is uniform throughout the cylinder and parallel to its unit axis. PhysLean's electric field retains the vector-field role; the dimensionful strength supplies its coherent-SI magnitude
\end{definition}
\begin{proof}
  This is the declared model definition of Is Uniform Electric Field Parallel To Cylinder Axis.
\end{proof}
\begin{definition}[axial Flux Orientation Factor]
  \label{def:physics:phyx-mini-0846:phyxminiproblems-problemphyxmini0846-axialfluxorientationfactor}
  \lean{PhyXMiniProblems.ProblemPhyXMini0846.axialFluxOrientationFactor}
  \uses{def:physics:phyx-mini-0846:phyxminiproblems-problemphyxmini0846-cylindersurfacepart}
  Orientation factor for the flux of a uniform axial field through each part of the closed cylinder: inward at the left cap, tangent to the wall, and outward at the right cap
\end{definition}
\begin{proof}
  This is the declared model definition of axial Flux Orientation Factor.
\end{proof}
\begin{definition}[Satisfies Closed Cylinder Electric Flux Laws]
  \label{def:physics:phyx-mini-0846:phyxminiproblems-problemphyxmini0846-satisfiesclosedcylinderelectricfluxlaws}
  \lean{PhyXMiniProblems.ProblemPhyXMini0846.SatisfiesClosedCylinderElectricFluxLaws}
  \uses{def:physics:phyx-mini-0846:phyxminiproblems-problemphyxmini0846-areainsquaremeters, def:physics:phyx-mini-0846:phyxminiproblems-problemphyxmini0846-electricfieldstrengthinnewtonspercoulomb, def:physics:phyx-mini-0846:phyxminiproblems-problemphyxmini0846-electricfluxinnewtonsquaremeterspercoulomb, def:physics:phyx-mini-0846:phyxminiproblems-problemphyxmini0846-uniformaxialfieldcylindersetup, def:physics:phyx-mini-0846:phyxminiproblems-problemphyxmini0846-axialfluxorientationfactor}
  Surface-flux evaluation and finite additivity for the closed cylinder. These are the governing flux-integral laws for a uniform field parallel to the axis; they state the three local contributions, not the requested net-zero answer
\end{definition}
\begin{proof}
  This is the declared model definition of Satisfies Closed Cylinder Electric Flux Laws.
\end{proof}
\begin{definition}[Answer Choice]
  \label{def:physics:phyx-mini-0846:phyxminiproblems-problemphyxmini0846-answerchoice}
  \lean{PhyXMiniProblems.ProblemPhyXMini0846.AnswerChoice}
  The answer labels printed beneath the problem
\end{definition}
\begin{proof}
  This is the declared model definition of Answer Choice.
\end{proof}
\begin{definition}[answer Choice Flux Readout]
  \label{def:physics:phyx-mini-0846:phyxminiproblems-problemphyxmini0846-answerchoicefluxreadout}
  \lean{PhyXMiniProblems.ProblemPhyXMini0846.answerChoiceFluxReadout}
  \uses{def:physics:phyx-mini-0846:phyxminiproblems-problemphyxmini0846-answerchoice}
  Numerical answer-choice readouts in N m² / C; -6/5 is the exact value represented by the printed decimal -1.2
\end{definition}
\begin{proof}
  This is the declared model definition of answer Choice Flux Readout.
\end{proof}
% --- Archon declaration topology end ---
% --- Archon physics formalization source end ---
